% ============================================================
% Section 2 -- Related Work
% Author (DG): Tran Tuan Khang (SE204881)
% Target file: paper/sections/02_related.tex
% ============================================================

\section{Related Work}
\label{sec:related}

Generating BDD scenarios with language models forces researchers to deal with two completely different problems: getting the meaning right and getting the syntax to parse. Most papers look at just one of these problems at a time. When testing LLMs on raw user stories, people usually report really high accuracy but they use data that does not represent real-world diversity. 

% ------------------------------------------------------------------
\subsection{Gherkin / BDD Generation Using LLMs}
\label{subsec:rw-gherkin}

Take Rathnayake et al.~\cite{rathnayake2026bddgeneration} for example. They ran GPT-4, Claude 3 Opus, and Gemini 1.5 Flash on 500 stories and saw a 91.16\% BERTScore for GPT-4. They pulled all those stories from just one company's internal backlog. That makes it impossible for anyone else to reproduce the experiment or know if the model would work for a different business domain.

Sometimes researchers do use open data. Fernandes et al.~\cite{fernandes2025gherkin} pulled test cases from public sources to compare seven models, including GPT-3.5 Turbo and DeepSeek R1. The problem is they only used ten test descriptions in total to get a 0.84 METEOR score for Gemini. Ten descriptions is simply not enough data to prove anything statistically. Karpurapu et al.~\cite{karpurapu2024bddautomation} did a bit better by testing about 50 stories with PaLM-2 and GPT-3.5-Turbo, but they used Gherkin-lint to validate the output which is a problem because a static linter is not a real test runner and it will pass code that has the right keywords even if the logic is broken and the actual test crashes. dos Santos et al.~\cite{dossantos2025bddtestgeneration} just looked at how much coverage ChatGPT had on 34 stories and didn't even run the code. Narvaez et al.~\cite{narvaez2025lawtogherkin} tested Claude 3.7 on food-safety laws, which is totally different from normal software requirements. No one is actually parsing the output in a real testing environment on a large dataset while also checking semantic embeddings at the same time.

% ------------------------------------------------------------------
\subsection{Test-Generation Pipelines and Domain Applications}
\label{subsec:rw-pipelines}

Other teams try to build whole testing pipelines using LLMs. Ferreira et al.~\cite{ferreira2025industrialcasestudy} set up GPT-4 Turbo at a car company to write scenarios. They reported zero syntax errors for 50 scenarios. Fonseca et al.~\cite{fonseca2025mobileacceptance} pushed JIRA issues from the BMW app through DeepSeek-R1 and got 93.3\% correct syntax. These tools clearly work well for the specific teams that built them. We have no idea if that 90\% success rate holds up when the tool is dropped into a banking app or a social network.

Looking at unit-test generation shows the same kind of roadblocks. Zhang et al.~\cite{zhang2024testbench} got GPT-4 to cover 92.51\% of lines in Java classes. They actually had to write a custom bug-fix script because GPT-3.5 was failing 97.84\% of the time on syntax. Kavuri~\cite{kavuri2022testscriptautomation} tested GPT-4 on mixed requirements. Chang and Shirazi~\cite{chang2025llmassessment} found GPT-4o-mini was pretty stable for Python code. Wang et al.~\cite{wang2024xuatcopilot} used three GPT agents to test WeChat Pay and got an 88.55\% pass rate on GUI scripts. Clicking buttons on a screen requires a completely different thought process than writing BDD scenarios.

% ------------------------------------------------------------------
\subsection{Research Gap and Positioning}
\label{subsec:rw-gap}

Looking at what everyone has done so far leaves a few obvious holes in the research. Nobody is using GPT-5.4-mini (API 2026-03-17) as the main engine to write Gherkin from Connextra stories. People rely on BERTScore or human grading for semantics instead of setting a hard 0.85 cosine cutoff with something like \texttt{all-MiniLM-L6-v2}. Almost everyone uses linters to check syntax rather than running the code in \texttt{behave --dry-run}. And there are no open multi-domain datasets with more than about 50 stories in this area.

Our experiment tries to fix these issues. We downloaded 261 Connextra stories from Apache Fineract, Diaspora, and Sylius. We ran them through GPT-5.4-mini. We gave the model one specific example for each project's domain. We checked the results against two hard rules: a 0.85 skeleton cosine threshold and an 80\% executable pass rate under \texttt{behave}. We haven't seen any other study run both of these strict checks on a dataset this big.
